\chapter{Diseño del sistema}
\label{anexo2:cap:diseno}
A partir del análisis del sistema, se puede proceder a su diseño. El diseño del sistema es una
representación más detallada de como se implementará el sistema, incluyendo la arquitectura,
los patrones de diseño utilizados, las tecnologías empleadas, etc., por lo que se puede considerar
una representación muy concreta de la solución planteada en el análisis \parencite{IBMDesignModel}.

Este capítulo se divide en varias secciones que describen los diferentes aspectos mencionados
anteriormente. En primer lugar, se presenta los patrónes de diseño utilizados para
construir el sistema. A continuación se muestra el diagrama de clases, que representa todos los objetos
y sus relacciones en el sistema. Finalmente, se presenta la realización de los casos de uso (definidos
en el \primerAnexo) y el diagrama de despliegue, que muestra como se desplegará
el sistema en un entorno real.

\section{Patrones de diseño}
\label{anexo2:sec:patron-diseno}
En el diseño del sistema se han utilizado varios patrones de diseño arquitectónico y de software. En
esta sección se describen los patrones utilizados y su justificación. Esto será de gran ayuda para
entender la estructura del sistema y el posterior diagrama de clases
(\ref{anexo2:fig:diagrama-clases-diseno}).

\subsection{Cliente-servidor}
\label{anexo2:sec:cliente-servidor}
El patrón de diseño arquitectónico escogido para el sistema es el patrón cliente-servidor. Este patrón
divide las tareas del sistema entre los proovedores de servicios (servidores) y aquellos que consumen
dichos servicios (clientes) \parencite{SunDistributedSystems}. En el caso de este sistema en concreto,
nos encontramos con un servidor que proporciona las funcionalidades necesarias para el funcionamiento
del sistema y otro servidor que proporciona la interfaz de usuario, que es consumida por los clientes
a través de un navegador web. Esta separación en tres capas, presentación, lógica de negocio y
acceso a datos, se conoce como arquitectura en tres capas \parencites{Ibm2025Apr}. Esta arquitectura
permite desacoplar la estructura de los datos de la presentación de los mismos, lo que facilita
la evolución del sistema y la reutilización de los componentes \parencite[p.~3]{nyabuto2024architectural}.

\subsection{\glsfirst{dao}}
\label{anexo2:sec:dao}
El servidor encargado de proporcionar las funcionalidades del sistema utiliza el patrón de diseño
\gls{dao} para acceder a los datos almacenados en la base de datos. Este patrón permite separar la lógica
de negocio del acceso a datos, lo que reduce el acoplamiento entre las diferentes capas del sistema y
facilita la modificación y mantenimiento de distintas secciones del sistema de manera independiente
\parencite{Zato2020Dao}.

% https://commons.wikimedia.org/wiki/File:IdeaDAOpng.png
\diagrama{
    nombre = {DAO (Imagen de flanciskinho, Creative Commons)},
    archivo = {anexo2/imgs/dao.png},
    etiqueta = {anexo2:fig:dao},
}

\section{Diagrama de clases}
\label{anexo2:sec:diagrama-clases-diseno}
El diagrama de clases de diseño, a diferencia del de análisis (\ref{anexo2:fig:modelo-dominio}), es
mucho más detallado, mostrando las clases tal y como se implementarán en el sistema, incluyendo
los atributos y métodos de cada clase, así como las relaciones entre ellas.

\section{Realización de los casos de uso}
\label{anexo2:sec:realizacion-casos-uso-diseno}

\section{Diagrama de despliegue}
\label{anexo2:sec:diagrama-despliegue}
